\documentclass{article}
\usepackage[T1]{fontenc}
\usepackage{lmodern}
\usepackage{graphicx}
\usepackage{amsmath,amssymb}
\usepackage[a4paper,margin=2cm]{geometry}
\usepackage[absolute,overlay]{textpos}
\setlength{\TPHorizModule}{1mm}
\setlength{\TPVertModule}{1mm}
\usepackage[indonesian]{babel}
\usepackage{hyperref}
\setlength{\emergencystretch}{2em}
% unit: Halaman 5

\begin{document}

% === Halaman 5 (llm) ===

\begin{itemize}
  \item $\langle G, + \rangle$ menyatakan sebuah grup dengan operasi penjumlahan.
  \item $\langle G, \cdot \rangle$ menyatakan sebuah grup dengan operasi perkalian
\end{itemize}

Contoh-contoh grup:

\begin{enumerate}
  \item $\langle R, + \rangle$ : grup dengan himpunan bilangan riil dengan operasi $+$ \\
  $e = 0$ dan $a' = -a$
  
  \item $\langle R^*, \cdot \rangle$ : grup dengan himpunan bilangan riil tidak nol (yaitu, $R^* = R - \{ 0 \}$) dengan operasi kali $(\cdot)$ \\
  $e = 1$ dan $a' = 1/a = a^{-1}$
  
  \item $\langle Z, + \rangle$ dan $\langle Z, \cdot \rangle$ masing-masing adalah grup dengan himpunan bilangan bulat (integer) dengan operasi $+$ dan $\cdot$.
\end{enumerate}

\vfill
\begin{center}
  Rinaldi Munir/II4021 Kriptografi
\end{center}

\end{document}
